\documentclass[11pt,a4paper]{article}
\usepackage[utf8]{inputenc}
\usepackage[T1]{fontenc}
\usepackage[french]{babel}
\usepackage{amsmath,amssymb,amsthm,booktabs,enumitem,geometry}
\geometry{margin=2.4cm}
\newtheorem{theoreme}{Théorème}
\newtheorem{proposition}[theoreme]{Proposition}
\newtheorem{lemme}[theoreme]{Lemme}
\newtheorem{corollaire}[theoreme]{Corollaire}
\theoremstyle{remark}
\newtheorem{remarque}[theoreme]{Remarque}
\newcommand{\R}{\mathbb R}
\newcommand{\norm}[1]{\lVert #1\rVert}
\newcommand{\vth}{\vartheta}
\newcommand{\D}{\mathrm D}
\sloppy\setlength{\emergencystretch}{3em}

\title{La criticité transverse est-elle une propriété structurelle\\ de certaines classes du noyau IWS ?}
\author{Programme de soustraction --- IWS, Tome II (travail de recherche)}
\date{octobre 2026 --- version V5}
\begin{document}
\maketitle

\noindent\textbf{Objet et statut.} La condition (U3) du complément « transmission $C^2$ » (criticité $\D a_\varepsilon=0$ sur le lieu des zéros du déplacement transverse) est la vraie hypothèse géométrique de la récurrence. On demande ici seulement : \emph{est-elle structurelle dans certaines classes du noyau, et quel mécanisme la favorise ou la détruit ?} Résultats : (i) un critère au premier ordre (section~\ref{sec:A}) ; (ii) une classe isotrope ($\lambda=0$, forçage axial) où (U3) est démontrée par invariance d'un sous-espace (section~\ref{sec:iso}) ; (iii) une classe anisotrope ($\lambda>0$) où (U3) \emph{échoue} au voisinage du sous-espace axial sous une condition de non-dégénérescence que la réponse linéaire calculée vérifie numériquement (section~\ref{sec:aniso}) ; (iv) un contre-exemple \emph{sans mot} dans le premier épisode (seuil de pression $\Theta=\frac{11}5$), \textbf{assisté par ordinateur} : (N) y est certifiée par arithmétique de boules, d'où la perte de $C^2$ de $P_0$ et de l'instant d'impact $T_1$ (sections~\ref{sec:premier}--\ref{sec:certif}) ; (v) le contact $a_\varepsilon\ge0$ n'est pas structurel dans ce cadre (section~\ref{sec:contact}). Les points (i)--(ii) sont analytiques ; (iv) est démontré modulo la certification numérique rigoureuse ; (iii) mêle un argument analytique (symétrie, formule de réponse) et des témoins numériques sur le noyau complet ; tout est \textbf{à relire}. Notation : $e=(1,0)$, $n=(0,1)$, $\vth=\tau(0)\in\R^{N\times2}$, $(\gamma_1,\lambda)=\varepsilon(\bar\gamma,\bar\lambda)$.

\section{Critère au premier ordre}\label{sec:A}
Soit un épisode de référence ($\varepsilon=0$, donc $W_0$ indépendante de $\vth$) avec un zéro simple $s_*$ : $W_0(s_*)=0$, $w=\partial_sW_0(s_*)\ne0$, $e_w=w/\lVert w\rVert$, $n_w\perp e_w$. À mot fixé, $(\varepsilon,\vth)\mapsto W_\varepsilon(s,\vth)$ est de classe $C^2$ (le champ dépend de $\varepsilon$ de façon régulière) : $W_\varepsilon=W_0+\varepsilon W_1+O(\varepsilon^2)$, uniformément en $C^1_\vth$. La projection fixe sur $e_w$ donne alors
\begin{equation}\label{eq:A}
a_\varepsilon(\vth)=\varepsilon A(\vth)+O(\varepsilon^2),\qquad A(\vth)=n_w\cdot W_1(s_*,\vth),
\end{equation}
car le déplacement $\sigma_\varepsilon-s_*=O(\varepsilon)$ du zéro temporel agit le long de $w$ et ne contribue pas à la composante normale au premier ordre.

\begin{proposition}[dichotomie au premier ordre]\label{prop:dich}
(a) Si $\lvert A\rvert\ge c>0$ sur la boule élargie, $a_\varepsilon$ ne s'annule pas pour $\varepsilon$ petit : (U3) est satisfaite \emph{vacuement}. (b) Si $A(\vth_*)=0$ et $\D A(\vth_*)\ne0$, alors $a_\varepsilon$ a un zéro régulier voisin de $\vth_*$ (théorème des fonctions implicites, uniformité $C^1$ de~\eqref{eq:A}) : (U3) échoue pour tout $\varepsilon$ petit, et la note « ordre » (nécessité) donne la perte de $C^2$ de l'intégrale de la norme. (c) Le cas intéressant est $A(\vth_*)=0$, $\D A(\vth_*)=0$.
\end{proposition}

\begin{remarque}[un contact au premier ordre ne suffit pas]\label{rem:contact1}
Si $a_\varepsilon(\vth)=\varepsilon\vth^2-\varepsilon^2$ : $A(\vth)=\vth^2$ (contact en $0$), mais $a_\varepsilon$ a deux zéros $\pm\sqrt\varepsilon$ avec $a_\varepsilon'=\pm2\varepsilon^{3/2}\ne0$ : (U3) échoue. Seule une identité exacte de signe (ou de carré) donnerait (U3) ; un contact observé à l'ordre dominant est insuffisant.
\end{remarque}

\section{Classe isotrope : $\lambda=0$, forçage axial}\label{sec:iso}
\begin{proposition}[invariance de l'axe]\label{prop:iso}
Prenons le noyau avec $\lambda=0$, $\gamma_1\ge0$ quelconque, forçage $u_i(t)\in\R e$ ($e$ un axe de \emph{coordonnée}), états initiaux $H_i(0),V_i(0)\in\R e$, et traces initiales $\vth=\tau(0)$ \emph{arbitraires}. Alors, pour tout $\vth$, tout épisode, $H_i(t),V_i(t)\in\R e$ ; en particulier $a_{\varepsilon,k}\equiv0$ pour tout zéro temporel de $H_i$ ou $V_i$.
\end{proposition}
\begin{proof}
Le sous-espace $E=\{H_i,V_i\in\R e\ \forall i\}$ est invariant par le champ à $\lambda=0$ : $M_\lambda(\tau_i)=I$ ; la friction $\gamma_0+\gamma_1\lVert\tau_i\rVert$ est scalaire ; $\widetilde A$ agit sur l'\emph{indice de nœud}, donc $(\widetilde AH)_i\in\R e$ si $H\in E$, et $\tanh$ s'applique coordonnée par coordonnée avec $\tanh(0)=0$ pour la composante $n$ ; le forçage est axial ; $\dot\tau$ n'intervient dans $(H,V)$ que par la norme $\lVert\tau_i\rVert$. Le reset $V_{i_j}\mapsto r_VV_{i_j}$ préserve $E$. Les instants et nœuds de Kairos (maximum de $P_i$, dynamique de $P$ par les normes $\lVert H_i\rVert=\lvert h_i\rvert$, $\lVert V_i\rVert=\lvert v_i\rvert$) et la réécriture du graphe (similarités cosinus des \emph{traces} seulement) ne dépendent pas de la direction de $H,V$ dans $E$. Par unicité des solutions d'un champ lipschitzien, la trajectoire issue de $E$ y reste. $\square$
\end{proof}
\begin{corollaire}[(U3) structurelle dans la classe isotrope]\label{cor:iso}
Dans la classe de la proposition~\ref{prop:iso}, avec zéros simples en $s$ des composantes scalaires $v_i,h_i$ ($i\in I_j$) et mot régulier, $Z=$ boule entière et $\D a\equiv0$ : (Z2) est vérifiée, aucun terme logarithmique n'apparaît ($T_1\equiv0$ dans l'identité du Hessien), et les théorèmes de transmission qualitative et quantitative du complément s'appliquent sous les seules hypothèses (D) et (U1). Les traces n'ont pas besoin d'être axiales : seules $H,V$ le sont.
\end{corollaire}
Note de portée : la classe est non générique (forçage et états initiaux axiaux), mais c'est une vraie sous-classe du noyau où les zéros de $V_0$ sont \emph{génériques} (changements de signe d'une fonction scalaire) et (U3) est garantie.

\section{Classe anisotrope : $\lambda>0$}\label{sec:aniso}
\paragraph{Symétrie.} Soit $R=\mathrm{diag}(1,-1)$ (réflexion $n\mapsto-n$), appliquée à tous les vecteurs $H_i,V_i,\tau_i$. Avec forçage axial, le champ du noyau est $R$-équivariant : $M_\lambda(R\tau)=RM_\lambda(\tau)R$ (car $\tau\tau^\top$ l'est), friction scalaire invariante, $\tanh$ coordonnée par coordonnée équivariante ; les similarités cosinus sont $R$-invariantes. Donc le flot est équivariant et, si $V_0=(V_0^x,V_0^n)$,
\[
V_0^x(s;R\vth)=V_0^x(s;\vth),\quad V_0^n(s;R\vth)=-V_0^n(s;\vth)\ \Longrightarrow\ \sigma_\varepsilon(R\vth)=\sigma_\varepsilon(\vth),\quad a_\varepsilon(R\vth)=-a_\varepsilon(\vth).
\]
Ainsi $a_\varepsilon=0$ sur $\mathrm{Fix}(R)=\{\text{traces axiales}\}$ pour tout $\varepsilon$ (zéros vectoriels persistants), et~(U3) en un point $p\in\mathrm{Fix}(R)$ équivaut à la nullité de la dérivée de $a_\varepsilon$ dans les directions \emph{transverses} $\vth^n$ (les dérivées le long de $\mathrm{Fix}(R)$ sont nulles).

\paragraph{Réponse transverse au premier ordre.} Posons $X_H=\partial_\varepsilon H^n_\varepsilon|_{\varepsilon=0}$, $X_V=\partial_\varepsilon V^n_\varepsilon|_{\varepsilon=0}$ (coefficients du premier ordre, sans facteur $\varepsilon$), calculés le long de la trajectoire de référence axiale ($\varepsilon=0$, qui ne dépend pas de $\vth$ pour $H,V$). Ils obéissent à un système \emph{linéaire à coefficients constants par épisode} :
\begin{equation}\label{eq:T}
\dot X_H=X_V,\qquad \dot X_V=-\gamma_0X_V-X_H+\zeta\widetilde A_kX_H+\bar\lambda f(s;\vth),\qquad f_i=\frac{h_i^0(s)\,\tau_i^x(s)\,\tau_i^n(s)}{1+(\tau_i^x)^2},
\end{equation}
avec $\tau_i^n(s)=e^{-\beta s}\vth_i^n$ (ordre dominant ; le terme en $(\tau^n_i)^2$ du dénominateur est d'ordre $2$ en $\vth^n$), $\tau_i^x$ la trace axiale de référence, $h^0$ la composante axiale de $H$ de référence ; sur plusieurs épisodes, $X_{V,i_j}\mapsto r_VX_{V,i_j}$ au Kairos et $\widetilde A_k\to\widetilde A_{k+1}$ (la dérivée de $\tanh$ en $0$ vaut $1$ pour la composante $n$ ; le terme $-\gamma_1\lVert\tau_i\rVert V^n$ est d'ordre $\varepsilon^2$ car $V^n=O(\varepsilon)$). Alors $a_\varepsilon=\varepsilon A+O(\varepsilon^2)$ avec $A(\vth)=X^{(\bar\lambda)}_{V,0}(s_*)$, où $X^{(\bar\lambda)}$ désigne la solution de~\eqref{eq:T} \emph{dont la source contient déjà $\bar\lambda$} (elle est linéaire en $\bar\lambda$ : $X^{(\bar\lambda)}=\bar\lambda X^{(1)}$, $X^{(1)}$ étant la solution pour $\bar\lambda=1$) (déplacement $\sigma_\varepsilon-s_*=O(\varepsilon)$ sans effet à cet ordre, car $V^n$ au zéro est $O(\varepsilon)$ et $\partial_sV^n=O(\varepsilon)$), et
\[
\D A(p)\cdot h=\bar\lambda\sum_i h_i^n\int_0^{s_*}G_{0i}(s_*,s)\,\frac{h_i^0(s)\,\tau_i^x(s)\,e^{-\beta s}}{1+\tau_i^x(s)^2}\,ds ,\qquad p\in\mathrm{Fix}(R),
\]
où $G$ est la fonction de Green (avec sauts) du système linéaire. C'est le seul mécanisme transverse : le terme de friction (tangentiel) ne contribue pas, d'où $A\equiv0$ pour $\bar\lambda=0$ (cohérent avec la proposition~\ref{prop:iso}).

\begin{proposition}[échec de (U3) sous non-dégénérescence]\label{prop:aniso}
Si $\bar\lambda>0$ et la forme linéaire $h\mapsto\D A(p)\cdot h$ est non nulle en un point $p\in\mathrm{Fix}(R)$ où un zéro simple de la référence existe (condition (N), à vérifier sur la formule ci-dessus), alors, \emph{sur une branche régulière fixée} (graphe et mot constants sur un voisinage ouvert de $p$ dans l'espace des traces, de sorte que l'expansion~\eqref{eq:A} soit valable en $C^1_\vth$), pour tout $\varepsilon$ petit $\D a_\varepsilon(p)\ne0$ alors que $a_\varepsilon(p)=0$ : $p$ est un zéro régulier de $a_\varepsilon$, (U3) échoue au voisinage de $p$, et l'intégrale $\int g\lVert V_0\rVert\,ds$ n'est $C^2$ sur aucun voisinage de $p$ (note « ordre », proposition de nécessité). En particulier un contact $a_\varepsilon\ge0$ est impossible près de $p$ : $a_\varepsilon$ est impaire en $\vth^n$ avec $\D_{\vth^n}a_\varepsilon(p)\ne0$.
\end{proposition}
\begin{proof}[Esquisse]
$a_\varepsilon(p)=0$ par la symétrie ; $\D a_\varepsilon(p)=\varepsilon\D A(p)+O(\varepsilon^2)\ne0$ par la proposition~\ref{prop:dich} et la formule de réponse. La conclusion sur l'intégrale est la proposition de nécessité (fenêtre autour de $s_*$, zéro simple). $\square$
\end{proof}
\noindent\emph{Conséquence sur l'instant d'impact (esquisse).} Si le nœud déclencheur est celui dont la norme s'annule, $P_{i_j}=P_{\rm reg}+P_{\rm sing}$ avec $P_{\rm reg}\in C^2$ et la seconde différence de $P_{\rm sing}$ le long de $h$ équivalente à $c\,\delta^2\ln\frac1\delta$, $c=\frac{2\kappa_2e^{-\kappa_3(T-s_*)}}{\lvert v_s\rvert}(\D a\cdot h)^2$ ; en inversant $P(T(\vth),\vth)=\theta$, la seconde différence de $T$ est $-\Delta^2P_{\rm sing}/\partial_tP+O(\delta^2)$ : \textbf{la dérivée seconde de $T_j$ diverge en $-\frac{c}{\partial_tP}\ln\frac1\delta$}. L'inversion rigoureuse (régularité de $\partial_tP$ en $\vth$) n'est pas écrite ici.

\section{Contre-exemple sans mot : le premier épisode}\label{sec:premier}
La réserve principale sur la section précédente est que le point axial est un mur du mot. On l'élimine en plaçant le zéro de $V_0$ \emph{avant le premier Kairos} : tant que $t<T_1$ le graphe est fixe ($A_1$ « kept1 »), aucune réécriture n'a eu lieu, et le système est une EDO à graphe constant dont le flot est analytique en $(\varepsilon,\vth)$ tant que $\tau_i\ne0$. Aucune similarité cosinus n'intervient avant l'impact.

\paragraph{Protocole $\Pi$.} Noyau complet, $N=6$, $d=2$, forçage constant $U=e$ sur le nœud~0, départ au repos, traces $\tau_i(0)=\tfrac35\,e$ (point $p\in\mathrm{Fix}(R)$), \textbf{pressions initiales $P_0(0)=r_P\Theta=\tfrac{77}{100}$ et $P_i(0)=0$ pour $i=1,\dots,5$}, $(\gamma_1,\lambda)=\varepsilon(\bar\gamma,\bar\lambda)$, $\bar\gamma\ge0$, $\bar\lambda>0$. Constantes du manuscrit ($\zeta=\frac45$, $\alpha=\frac9{20}$, $\beta=\frac3{25}$, $\gamma_0=\frac14$, $\kappa_1=\frac35$, $\kappa_2=\frac25$, $\kappa_3=\frac12$, $r_P=\frac7{20}$), \textbf{sauf le seuil de pression $\Theta=\frac{11}5$ au lieu de $\frac54$.} Avec $\Theta=1{,}25$ le premier impact a lieu à $T_1=2{,}38$, avant tout zéro de $V_0^x$ ; il faut $\Theta\gtrsim2{,}0$ pour que le premier extremum de $H_0^x$ (zéro de $V_0^x$, indépendant de $\Theta$) précède l'impact. Le seuil $\Theta>0$ est un paramètre de la famille du noyau (Tome I, section « Ensemble de flot et ensemble de saut ») : $\Theta=\frac{11}5$ appartient à la famille, même s'il diffère du protocole numérique de référence $\Theta=\frac54$. La pression initiale est une hypothèse du protocole, non un détail : elle ne rétroagit pas sur $(H,V,\tau)$, donc ne change ni le zéro $s_*$ ni les données transverses, mais elle fixe l'instant d'impact (à $\varepsilon=0$, $T_1=4{,}455195$ avec $P_0(0)=\frac{77}{100}$, $T_1=4{,}812400$ avec $P_0(0)=0$).

\paragraph{Données de référence ($\varepsilon=0$), certifiées (section~\ref{sec:certif}).} Zéro simple unique $s_*^0=3{,}7013376\ldots$ de $v_0:=V_{0}^x$ sur $]0,4{,}6]$, $\partial_sv_0(s_*^0)=-0{,}5578391\ldots<0$ ; $v_0>0$ avant, $v_0<0$ après ; $\tau_i^x\ge0{,}46$ pour tout $i$ ; $V_0=s\,\tilde v_0$, $H_0=s^2\tilde h_0$ avec $\tilde v_0\ge0{,}993$, $\tilde h_0\ge0{,}497$ sur $]0,\frac1{20}]$ ; $h_0\ge0{,}117$ sur $[\frac12,4{,}6]$ ; $P_i\le2{,}131<\Theta$ sur $[0,4]$ pour tout $i$, $P_i\le0{,}74$ ($i\ge1$) sur $[0,4{,}6]$, $P_0$ strictement croissante tant qu'elle est $\le\Theta$ sur $[4,4{,}6]$ et $T_1^0\in\,]4{,}45;4{,}50[$, avec $\partial_tP_0(T_1^0)\ge0{,}214$ ; le déclencheur est le nœud~0 (strictement). Par continuité du flot en $(\varepsilon,\vth)$ et fonctions implicites, il existe $\varepsilon_0>0$ (non explicité) tel que pour $0<\varepsilon<\varepsilon_0$ et $\vth$ voisin de $p$, le zéro $s_*^\varepsilon(\vth)$ de $V_0^x$ est simple, de classe $C^\infty$, et ces propriétés persistent ; on fixe $[t_1,t_2]=[3{,}75;\,4{,}55]$, avec $s_*^0<t_1$ et $T_1^0<t_2$. Au-delà de l'impact on utilise la \emph{prolongation virtuelle} du flot (EDO sans saut), notée encore $P_0(t;\vth)$, qui coïncide avec la pression réelle pour $t<T_1(\vth)$. L'analyticité en $(t,\varepsilon,\vth)$ est celle du sous-système $(H,V,\tau)$ ; la pression prolongée est continue (et dérivable en $\vth$), mais précisément pas de classe $C^2$ en $\vth$ au zéro transverse, ce que le théorème ci-dessous établit : c'est à elle que s'applique le lemme logarithmique. (À $\varepsilon=0{,}1$, $s_*^\varepsilon=3{,}836>t_1$ : $\varepsilon_0$ est donc plus petit que $0{,}1$, mais $t=4{,}0$ reste admissible pour les tests numériques.)

\begin{theoreme}[contre-exemple assisté par ordinateur : perte de $C^2$ dans le premier épisode]\label{prop:premier}
Sous le protocole $\Pi$, pour $0<\varepsilon<\varepsilon_0$ fixé et tout $t\in[t_1,t_2]$ (pression prolongée virtuellement au-delà de l'impact) : (i) $a_\varepsilon(p)=0$ et $\D a_\varepsilon(p)\cdot e_{0n}=\varepsilon\bar\lambda\cdot0{,}4954739718\ldots+O(\varepsilon^2)\ne0$ ($e_{0n}$ : direction normale de la trace du nœud~0) ; (ii) $\vth\mapsto P_0(t;\vth)$ n'est pas de classe $C^2$ au voisinage de $p$ ; (iii) pour $h=e_{0n}$,
\begin{gather*}
P_0(t;p+\delta h)+P_0(t;p-\delta h)-2P_0(t;p)=c_\varepsilon(t)\,\delta^2\ln\tfrac1{\lvert\delta\rvert}+O(\delta^2),\\
c_\varepsilon(t)=\frac{2\kappa_2\,e^{-\kappa_3\,(t-s_*^\varepsilon)}}{\bigl\lvert\partial_sV^x_{0,\varepsilon}(s_*^\varepsilon,p)\bigr\rvert}\,\bigl(\D a_\varepsilon(p)\cdot h\bigr)^2>0 ,
\end{gather*}
où $s_*^\varepsilon=\sigma_\varepsilon(p)$ est le zéro \emph{perturbé} de $V^x_{0,\varepsilon}(\cdot,p)$ (et non $s_*^0$) : le coefficient $c_\varepsilon$ s'évalue à $\varepsilon$ fixé, ses valeurs en $\varepsilon=0$ n'en sont qu'une approximation perturbative.
\end{theoreme}
\begin{proof}
\emph{Régularité avant le zéro.} Avant $T_1$ aucun saut n'a lieu : la branche est triviale, $(s,\varepsilon,\vth)\mapsto(H,V,\tau)$ est analytique, et le zéro simple $\sigma_\varepsilon(\vth)$ de $V_0^x$ l'est aussi (fonctions implicites, $\partial_sv_0\ne0$). Le développement~\eqref{eq:A} vaut en $C^1_\vth$ et $A(\vth)=X^{(\bar\lambda)}_{V,0}(s_*^0)$ pour la solution de~\eqref{eq:T}. \emph{(i)} $a_\varepsilon(p)=0$ par la symétrie $R$ (section~\ref{sec:aniso}) ; la dérivée vaut $\D A(p)\cdot e_{0n}=X^{(\bar\lambda)}_{V,0}(s_*^0)=\bar\lambda\,X^{(1)}_{V,0}(s_*^0)$ (le facteur $\bar\lambda$ n'est compté qu'une fois : il est dans la source, et $X^{(1)}$ est obtenue en dérivant~\eqref{eq:T} en $\vth^n_0$ pour $\bar\lambda=1$ ; seule la source du nœud~0 survit, $\tau_0^n=e^{-\beta s}$, dénominateur $1+(\tau_0^x)^2$), soit $0{,}49547397\ldots\,\bar\lambda\ne0$ par le théorème~\ref{thm:certif} ; d'où $\D a_\varepsilon(p)\cdot e_{0n}=\varepsilon\D A(p)\cdot e_{0n}+O(\varepsilon^2)\ne0$.

\emph{(ii)--(iii).} $P_0(t;\vth)=P_0(0)\,e^{-\kappa_3t}+\kappa_1\!\int_0^t\!e^{-\kappa_3(t-s)}\lVert H_0\rVert\,ds+\kappa_2\!\int_0^t\!e^{-\kappa_3(t-s)}\lVert V_0\rVert\,ds$, avec $P_0(0)=\frac{77}{100}$ : le premier terme est indépendant de $\vth$ (donc régulier, et sans effet sur le coefficient logarithmique) mais il est indispensable pour retrouver l'instant d'impact $T_1$. \emph{Traitement de $s=0$.} Le départ au repos donne $V_0(s)=s\,\tilde v_0(s)$, $H_0(s)=s^2\tilde h_0(s)$ avec $\tilde v_0,\tilde h_0$ lisses en $(s,\varepsilon,\vth)$ et $\lVert\tilde v_0\rVert\ge0{,}993$, $\lVert\tilde h_0\rVert\ge0{,}497$ sur $[0,\frac1{20}]$ (les minorations de $\lVert H_0\rVert$ sur $[\frac1{20},4{,}6]$ ne traitent pas cette extrémité) : sur $[0,\frac1{20}]$ les intégrandes $s\lVert\tilde v_0\rVert$, $s^2\lVert\tilde h_0\rVert$ sont lisses, et les contributions correspondantes sont $C^\infty$ en $\vth$. \emph{Reste.} Sur $[\frac1{20},t]$ : $\lVert H_0\rVert=h_0$ est lisse et minoré par la borne certifiée $h_0\ge0{,}000626$ sur $[\frac1{20},4{,}6]$ (théorème~\ref{thm:certif}), et même $h_0\ge0{,}117$ sur $[\frac12,4{,}6]$ (la factorisation ne sert que sur $[0,\frac1{20}]$) ; ces minorations persistent pour $\varepsilon$ et $\vth-p$ petits par continuité ; $V_0(\frac1{20})\ne0$ et $V_0(t)\ne0$ (aux deux bornes, $v_0>0$ puis $v_0<0$), $V_0(\cdot,p)$ a un unique zéro simple $s_*^\varepsilon$ intérieur à $[\frac1{20},t]$, de composante transverse $a_\varepsilon(\vth)$ de gradient $\D a_\varepsilon(p)\ne0$ : ce sont les hypothèses (H1)--(H3) du lemme logarithmique, avec poids $g(s)=\kappa_2e^{-\kappa_3(t-s)}$ indépendant de $\vth$. Il donne l'estimation de la différence seconde, avec reste $O(\delta^2)$ (formule de la différence seconde de la note « singularité logarithmique »), donc (iii), et la perte de $C^2$ (ii) puisque $c_\varepsilon(t)>0$. $\square$
\end{proof}

\begin{proposition}[perte de $C^2$ de l'instant d'impact $T_1$]\label{cor:T1}
Sous le protocole $\Pi$, pour $0<\varepsilon<\varepsilon_0$ fixé, $t_0:=T_1^\varepsilon(p)$ et $h=e_{0n}$ :
\[
T_1(p+\delta h)+T_1(p-\delta h)-2T_1(p)=-\frac{c_\varepsilon(t_0)}{\partial_tP_0(t_0,p)}\,\delta^2\ln\frac1{\lvert\delta\rvert}+O(\delta^2),
\]
donc $T_1$ n'est pas de classe $C^2$ en $p$ (elle n'est pas deux fois dérivable en ce point).
\end{proposition}
\begin{proof}
Par la symétrie $R$, $P_0$ et $T_1$ sont pairs en $\delta$ ($p-\delta h=R(p+\delta h)$). Posons $F(\delta)=P_0(t_0,p+\delta h)$ : par le théorème~\ref{prop:premier} appliqué en $t=t_0\in[t_1,t_2]$ ($T_1^\varepsilon(p)$ est voisin de $T_1^0\in\,]4{,}45;4{,}50[$, donc dans l'intervalle ; $P_0(t_0,\cdot)$ est la pression prolongée virtuellement, égale à la pression réelle jusqu'à l'impact), $F(\delta)-F(0)=\tfrac12c_\varepsilon(t_0)\delta^2\ln\frac1{\lvert\delta\rvert}+O(\delta^2)$. Soit $\Delta(\delta)=T_1(p+\delta h)-t_0$. Par définition $P_0(t_0+\Delta,p+\delta h)=\Theta=P_0(t_0,p)$, d'où
\[
\Delta\,\partial_tP_0(\xi,p+\delta h)=-\bigl(F(\delta)-F(0)\bigr),\qquad\xi\text{ entre }t_0\text{ et }t_0+\Delta .
\]
La dérivée $\partial_tP_0=\kappa_1\lVert H_0\rVert+\kappa_2\lVert V_0\rVert-\kappa_3P_0$ est continue en $(t,\vth)$ et, au voisinage de $(t_0,p)$, où $\lVert H_0\rVert,\lVert V_0\rVert$ sont lisses (aucun zéro), sa variation vient de $P_0$ : $\partial_tP_0(\xi,p+\delta h)=\partial_tP_0(t_0,p)+O(\lvert\Delta\rvert+\delta^2\ln\frac1{\lvert\delta\rvert})$. Comme $\partial_tP_0(t_0,p)>0$ et $\Delta=O(\delta^2\ln\frac1{\lvert\delta\rvert})$,
\[
\Delta(\delta)=-\frac{c_\varepsilon(t_0)}{2\partial_tP_0(t_0,p)}\delta^2\ln\frac1{\lvert\delta\rvert}+O(\delta^2),
\]
le reste d'inversion étant d'ordre $\delta^4\ln^2\frac1{\lvert\delta\rvert}$ (et non $O(\delta^4)$). Par parité, la différence seconde de $T_1$ vaut $2\Delta(\delta)$. $\square$
\end{proof}

\paragraph{Témoins numériques} (\texttt{verif\_premier\_episode.py}, sortie \texttt{sortie\_premier\_episode.txt}).
\begin{enumerate}[nosep]
\item \emph{Réponse transverse, tous les nœuds valides} (aucun mot en jeu ; $\varepsilon=d=10^{-2}$) : pour $(\bar\gamma,\bar\lambda)=(1,0)$ : $0$ exactement ; pour $(0,1)$ : $0{,}4945$ (nœud~0), puis $1{,}26\cdot10^{-2}$, $1{,}2\cdot10^{-5}$, $3{,}3\cdot10^{-5}$, $7{,}95\cdot10^{-3}$, $7{,}14\cdot10^{-3}$ (nœuds $1,\dots,5$) ; pour $(1,1)$ : $0{,}4886$ ($\lVert\D A(p)\rVert=0{,}489$). La formule linéaire~\eqref{eq:T} donne $0{,}49547$ pour le nœud~0 (écart $0{,}21\%$) et reste à $1{,}6$--$2{,}8\%$ du noyau pour les cinq autres (compatible avec $O(\varepsilon)$). L'extrapolation de Richardson du noyau en $\varepsilon\to0$ (\texttt{verif\_N\_extrapolation\_eps.py} : $\varepsilon=10^{-2},5\cdot10^{-3},2{,}5\cdot10^{-3},1{,}25\cdot10^{-3}$) donne $0{,}49547385$ (ordre~2), contre l'encadrement certifié $0{,}4954739718\pm2\cdot10^{-13}$ : accord à $2\cdot10^{-7}$ en relatif.
\item \emph{Douze directions du sous-espace caché} : une perturbation axiale aléatoire des six traces donne $a=0$ ; pour une direction normale aléatoire $h\in\R^6$, $(a(+h)-a(-h))/(2\varepsilon)=2{,}60620\cdot10^{-4}$ contre $\sum_j\D A_jh_j=2{,}60605\cdot10^{-4}$ : la différentielle est bien la forme linéaire de coefficients ci-dessus.
\item \emph{Seconde différence} ($\varepsilon=0{,}1$, $\bar\gamma=1$, direction $n$ du nœud~0, $t=4{,}0$) : pour $\lambda=0$, $\D^2P_0$ et $\D^2T_1$ restent bornés ($-0{,}144538$ et $1{,}423247$ pour $\delta=10^{-3}$) ; pour $\lambda=\varepsilon$, pentes par décade de $\D^2P_0(t=4)$ en $\ln\frac1\delta$ : $+6{,}8544\cdot10^{-3}$ puis $+6{,}8698\cdot10^{-3}$ contre $+6{,}8714\cdot10^{-3}$ prédit (écarts $0{,}25\%$ et $0{,}02\%$) ; pour $\D^2T_1$, avec $T_1$ obtenu par racine de $P_0(t)=\Theta$ sur intégration coupée autour du zéro : $-2{,}3824\cdot10^{-2}$ puis $-2{,}3867\cdot10^{-2}$ contre $-2{,}3872\cdot10^{-2}$ prédit ($0{,}2\%$ et $0{,}02\%$). Les paramètres de la prédiction, \emph{évalués à $\varepsilon=0{,}1$} ($s_*^\varepsilon=3{,}8359$, $\partial_sV^x=-0{,}4517$, $a'=4{,}277\cdot10^{-2}$, $T_1=4{,}5918$, $\partial_tP_0=0{,}2141$), sont mesurés séparément, sans ajustement : c'est bien $c_\varepsilon$ et non $c_0$ qui est testé.
\item \emph{Sur le bruit de l'épisode~2} : la même seconde différence de $T_1$ calculée par localisation d'événement de \texttt{solve\_ivp} dérive comme dans l'épisode~2 ($-2{,}39\cdot10^{-2}$ puis $-2{,}60\cdot10^{-2}$ par décade), alors que le calcul par racine est stable. La dérive aux petits $\delta$ du point~4 de la section~\ref{sec:numerique} est donc très probablement un artefact de localisation d'événement et non une correction réelle ; je n'ai pas refait l'épisode~2 par la méthode de la racine, et ses pentes restent à lire avec cette réserve.
\end{enumerate}

\section{Certification de (N) par arithmétique de boules}\label{sec:certif}
\paragraph{Ce qui est certifié.} La quantité $X^{(1)}_{V,0}(s_*^0)$, qui donne $\D A(p)\cdot e_{0n}=\bar\lambda X^{(1)}_{V,0}(s_*^0)$, ne dépend que de la trajectoire de référence ($\varepsilon=0$, axiale, $H,V,\tau\in\R^{6}$ en composantes $e$) et du système linéaire~\eqref{eq:T} : c'est la valeur finale d'une EDO \emph{polynomiale} après introduction des variables $w_i=\tanh((\widetilde Ah)_i)$ ($\dot w_i=(1-w_i^2)(\widetilde Av)_i$), $E=e^{-\beta s}$ et $q=(1+\tau_0^2)^{-1}$, avec données initiales et coefficients rationnels exacts. L'intégrateur est un schéma de Taylor validé d'ordre $22$ en arithmétique de boules ($200$ bits, bibliothèque \texttt{python-flint}/Arb), pas $\frac1{20}$, avec enclosure a priori de Picard ($x_0+[0,h]f(B)\subset B$) et reste de Lagrange évalué sur $B$ ; le programme est \texttt{certif\_N\_premier\_episode\_V4.py} (sortie \texttt{sortie\_certification\_N\_V4.txt}, exécution en moins d'une seconde). Toutes les conclusions du théorème~\ref{thm:certif} sont des \emph{assertions} : le programme lève une \texttt{AssertionError} si l'une des inégalités n'est pas certaine (comparaisons entre boules Arb elles-mêmes, jamais entre flottants arrondis) ; les bornes décimales affichées sont arrondies vers l'extérieur ; une altération volontaire d'un seuil fait échouer l'exécution.

\begin{theoreme}[encadrements certifiés, protocole $\Pi$, $\varepsilon=0$]\label{thm:certif}
\begin{enumerate}[nosep]
\item $v_0:=V_0^x>0$ sur $]0,s_*^0[$ ; $v_0$ a un unique zéro $s_*^0\in[3{,}70133761238468\pm3\cdot10^{-15}]$ dans $]0,4{,}6]$ ; $\partial_sv_0\in[-0{,}5578391252\pm4\cdot10^{-11}]$ en $s_*^0$ (et $\partial_sv_0<-0{,}546$ sur $[3{,}70;3{,}75]$) ; $v_0<0$ sur $]s_*^0,4{,}6]$.
\item $\boxed{\;0{,}49547397\le X^{(1)}_{V,0}(s_*^0)\le0{,}49547398\;}$ (valeur centrale $0{,}495473971795$, rayon $2{,}1\cdot10^{-13}$ ; bornes décimales arrondies vers l'extérieur : $0{,}495473971794\le X^{(1)}_{V,0}(s_*^0)\le0{,}495473971795$), pour $\bar\lambda=1$ ; en particulier la forme linéaire $\D A(p)\cdot e_{0n}$ est non nulle, et plus précisément comprise dans $]0{,}49;0{,}50[$ pour $\bar\lambda=1$.
\item $\tau_i^x\ge0{,}46$ ($i=0,\dots,5$) sur $[0,4{,}6]$ ; $v_0/s\ge0{,}993$, $h_0/s^2\ge0{,}497$ sur $]0,\frac1{20}]$ ; $h_0\ge0{,}000626$ sur $[\frac1{20},4{,}6]$ et $h_0\ge0{,}117$ sur $[\frac12,4{,}6]$.
\item (avec $P_0(0)=\frac{77}{100}$, $P_i(0)=0$ pour $i\ge1$) $P_i(t)\le2{,}131<\Theta$ pour $t\le4$, tout $i$ ; $P_i\le0{,}74$ pour $i\ge1$, $t\le4{,}6$ ; sur $[4;4{,}6]$, $g=\kappa_1h_0+\kappa_2\lvert v_0\rvert$ vérifie $g-\kappa_3\Theta\ge0{,}214$ ; $P_0(4{,}45)\le2{,}19879<\Theta<2{,}21029\le P_0(4{,}50)$. Donc $T_1^0\in\,]4{,}45;4{,}50[$, simple, le nœud~0 déclenche seul.
\end{enumerate}
\end{theoreme}
Le lien avec le texte analytique est le suivant : le théorème~\ref{thm:certif} ne contient aucune hypothèse sur $\varepsilon>0$ ; le passage à $\varepsilon\in]0,\varepsilon_0[$ est la persistance analytique (fonctions implicites, théorème~\ref{prop:premier}). Les enclosures sont des \emph{preuves assistées par ordinateur}, correctes dès lors que le code et la bibliothèque d'arithmétique de boules le sont ; je ne les ai pas vérifiées formellement. Elles ont été recoupées par le solveur flottant (valeurs de $h_0,v_0,\tau_0,P_0,P_1$ à $t=1,2,3,4$ égales à $10^{-9}$ près à \texttt{solve\_ivp}) et par la formule indépendante du noyau complet (point~1 des témoins). Les données mathématiques du certificat (coefficients du noyau, matrice $\widetilde A$, conditions initiales, seuils) sont des rationnels exacts, encadrés ensuite par Arb ; toutes les décisions (signes, inégalités, choix du pas de pression) reposent sur des comparaisons entre boules Arb, sur des tubes de Taylor. Les conversions en flottants ne servent qu'à l'affichage. La bissection du zéro manipule des bornes rationnelles dont le signe est certifié par la boule.

\paragraph{Reproductibilité.} Environnement : Python $3.11.15$, \texttt{python-flint} $0.9.0$ (\texttt{pip install python-flint==0.9.0}), précision $200$ bits. Empreintes SHA-256 :
\begin{center}\footnotesize
\texttt{4125cf7426fdc5a4416743cf951b52e88b1e10722ecbea3580297e5ef8cf92eb}\ \ (\texttt{certif\_N\_premier\_episode\_V4.py})\\
\texttt{97b643bb1fa291f90f1acb5bc51b9be96400657b5019fbd3f6411d0f53e4d937}\ \ (\texttt{sortie\_certification\_N\_V4.txt})
\end{center}
Le fichier \texttt{MANIFEST\_certification\_V4.txt} les regroupe. La fonction de décision de signe utilise désormais directement les comparaisons Arb (depuis la version~V3 du script ; la V2 passait par des flottants et pouvait seulement recourir inutilement à la borne grossière). La version~V4 du script supprime un test redondant dans l'enclosure a priori (la boule $\mathrm{newB}=x_0+[0,h]f(B)\subseteq B$ est elle-même une enclosure valide de la solution sur $[0,h]$, par l'isotonie de l'évaluation par intervalles : le second test ne servait à rien) et une expression morte \texttt{if False else} dans le calcul de pression sur le pas du zéro ; la sortie est \emph{identique} octet pour octet à celle de la V3.

\paragraph{Portée du résultat.} Pour ce protocole $\Pi$ précis (seuil $\Theta=\frac{11}5$, traces initiales $\frac35e$, constantes du manuscrit), le théorème~\ref{prop:premier} et la proposition~\ref{cor:T1} donnent un \emph{contre-exemple à (U3) et à la régularité $C^2$ de $T_1$ dans le noyau lui-même}, entièrement établi modulo la certification assistée par ordinateur de~(N). Cela ne dit rien sur le protocole du manuscrit ($\Theta=1{,}25$), dont le premier épisode n'a pas de zéro ; ni sur tous les systèmes anisotropes ; en revanche cela montre que (U3) n'est pas une conséquence automatique de l'anisotropie faiblement couplée et qu'elle ne peut pas être retirée des hypothèses générales du noyau.

\section{Le contact non négatif}\label{sec:contact}
Une classe à contact exigerait $a_\varepsilon\ge0$ exactement (remarque~\ref{rem:contact1}). Dans le noyau, l'unique symétrie trouvée est $R$, qui rend $a_\varepsilon$ \emph{impaire} : la positivité forcerait donc $\D a_\varepsilon=0$ sur $\mathrm{Fix}(R)$, c'est-à-dire l'annulation de la réponse linéaire $\D A$ --- une condition de dégénérescence supplémentaire (la forme linéaire $\D A(p)$ vit sur toutes les directions transverses des traces et doit s'annuler identiquement), non imposée par la symétrie, donc non structurelle. Je n'ai pas trouvé de mécanisme imposant un signe ; en l'état, le contact n'est pas une classe naturelle du noyau.

\section{Contrôles numériques sur le noyau complet (épisode~2)}\label{sec:numerique}
Scripts : \texttt{crit\_lib.py}, \texttt{verif\_criticite\_classes.py} et \texttt{verif\_reponse\_transverse\allowbreak\_formule.py} (dossier \texttt{variation\_y\_trace\_nulle}) ; $N=6$, $d=2$, forçage constant $U=(1,0)$ sur le nœud~0, zéro simple de $V_0^x$ dans l'épisode~2 (point de base $\tau_i(0)=0{,}6\,e$ : $s_*=0{,}978$, $v_s=-0{,}331$ en $\varepsilon=0$ ; à $\varepsilon=0{,}1$ du point~4, $s_*=1{,}101$, $v_s=-0{,}2475$). Ce sont des témoins, pas des preuves.
\begin{enumerate}[nosep]
\item \emph{Invariance de l'axe} ($\lambda=0$, $\gamma_1=0{,}3$, 8 tirages $\vth$ aléatoires sur les 12 composantes, épisodes 1--3) : $\max\lvert H^n,V^n\rvert=0$ ; zéro simple de l'épisode 2 dans $8/8$ cas.
\item \emph{Réponse transverse} ($\varepsilon=10^{-2}$, perturbation $n$ de la trace du nœud $j$ de $\pm10^{-2}$, point de base $\tau_i(0)=0{,}6\,e$) : pour $(\bar\gamma,\bar\lambda)=(1,0)$ : $0$ exactement ; pour $(0,1)$ : nœud~0 $=0{,}3707$ (autres nœuds $\lesssim5\cdot10^{-4}$) ; pour $(1,1)$ : $0{,}3670$. La friction ne change pas l'ordre dominant. Perturbation axiale : $a=0$ ; imparité $a(+0{,}3)=-a(-0{,}3)$ exacte. Au point $\tau(0)=0$ : $0{,}2482$ (nœud~0).
\item \emph{Formule linéaire indépendante} ($\bar\gamma=0$) : intégration de~\eqref{eq:T} le long de la référence axiale : nœud~0 : $0{,}3699$ contre $0{,}3707$ (écart relatif $2\cdot10^{-3}$) ; nœuds~2,3,4 : écarts de $2$ à $4\%$ (compatible avec $O(\varepsilon)$). Pour les nœuds~1 et~5 la perturbation $n$ \emph{change le mot} (la réécriture du graphe par similarité cosinus est départagée par l'arrondi, voir ci-dessous) : la comparaison n'a pas de sens pour ces deux nœuds (écart de $17\%$ pour le nœud~5, facteur $27$ pour le nœud~1), ce que signale l'indicateur « mot conservé » de la sortie ; ils sont exclus du constat.
\item \emph{Instant d'impact} ($\varepsilon=0{,}1$, $\bar\gamma=1$, direction $n$ du nœud~0, $\delta\in[3\cdot10^{-4},10^{-1}]$, seconde différence de $T_2$) : pour $\lambda=0$, $\D^2T_2\to0{,}25109$ (borné, $C^2$) ; pour $\lambda=0{,}1$, $\D^2T_2$ décroît linéairement en $\ln\frac1\delta$, avec des pentes mesurées par décade $-1{,}47\cdot10^{-2},\,-1{,}72\cdot10^{-2},\,-1{,}742\cdot10^{-2},\,-1{,}83\cdot10^{-2},\,-2{,}57\cdot10^{-2}$ ; \emph{pente prédite} $-1{,}742\cdot10^{-2}$ par la formule $-c/\partial_tP$ ci-dessus ($v_s=-0{,}2475$, $a'=3{,}405\cdot10^{-2}$, $P_t=0{,}4674$). Il faut lire cet accord avec mesure : la pente est à $1{,}5\%$ de la prédiction sur $3\cdot10^{-2}\to10^{-2}$, coïncide à cinq chiffres sur $10^{-2}\to3\cdot10^{-3}$ (coïncidence de milieu de plage), puis dérive ($-1{,}83$ puis $-2{,}57\cdot10^{-2}$) aux petits $\delta$, où le bruit de la seconde différence (division par $\delta^2$) n'est pas séparé d'une éventuelle correction réelle ; au grand $\delta$ l'écart ($15\%$) est celui des termes en $\delta^2$. Le témoin établit donc une divergence logarithmique de pente de l'ordre prédit (à quelques pour cent sur la plage fiable), non une validation à cinq chiffres. La constante $K=c/\partial_tP$ est bien recalculée à partir des paramètres indiqués ($7{,}566\cdot10^{-3}$).
\end{enumerate}
\paragraph{Deux réserves de méthode.} (a) \emph{Le point de base axial est un « mur » du mot} : les traces étant colinéaires, toutes les similarités cosinus valent $1$ et le choix des arêtes retirées/ajoutées à la réécriture du graphe est décidé par l'arrondi flottant ; la famille testée (perturbation $n$ du nœud~0) conserve les égalités et donc le mot, mais le mot n'est pas \emph{robuste} en ce point. Pour la classe isotrope (traces génériques) les écarts sont stricts et le mot est régulier. Pour la classe anisotrope, le lieu des zéros non critiques s'étend à des traces légèrement non axiales ($\lvert A\rvert\to0$ avec $\vth^n$ petit, écarts de similarité petits mais stricts), mais la persistance du mot près de $\mathrm{Fix}(R)$ est une vérification supplémentaire non faite (condition (U1)). Cette réserve ne touche pas la section~\ref{sec:premier}, où aucun saut n'a lieu avant le zéro. (b) \emph{Point $\tau(0)=0$} : une première série de mesures de $\D^2T$ a été faite en $\tau(0)=0$ ; elle divergeait en $\delta^{-1/2}$ même pour $\lambda=0$ : la norme $\lVert\tau_i\rVert$ de la friction n'est pas différentiable en $\tau=0$ (traces nulles à $t=0$), ce qui viole l'hypothèse de traces non nulles de P2 et n'a rien à voir avec le zéro de $V_0$. D'où le point de base $\tau_i(0)=0{,}6\,e$.

\section{Où en est-on}
\begin{center}\small
\begin{tabular}{p{6.3cm}p{4.2cm}p{4.0cm}}
\toprule
Énoncé & Hypothèses & Statut\\\midrule
Dichotomie au premier ordre (prop.~\ref{prop:dich}) & expansion $C^1_\vth$ & Analytique, à relire\\
Invariance de l'axe, $a\equiv0$ pour $\lambda=0$ (prop.~\ref{prop:iso}) & forçage axial, $e$ axe de coordonnée & Démontré (invariance) ; témoin numérique\\
(U3) structurelle en classe isotrope (cor.~\ref{cor:iso}) & + zéros simples, mot régulier & Démontré, à relire\\
Symétrie $R$ : $a_\varepsilon$ impaire & forçage axial, $\lambda\ge0$ & Démontré (équivariance)\\
Échec de (U3) pour $\lambda>0$ (prop.~\ref{prop:aniso}) & (N) : $\D A(p)\ne0$ & Conditionnel à (N) ; (N) vérifiée numériquement (nœud~0 : $0{,}37$) ; formule indépendante à $2\cdot10^{-3}$\\
$T_j\notin C^2$ pour $\lambda>0$, épisodes $\ge2$ & + inversion, mot & Esquissé ; divergence log observée, pente prédite à quelques \% sur la plage fiable\\
(N), premier épisode, $\Theta=\frac{11}5$ (thm.~\ref{thm:certif}) & arithmétique de boules & \textbf{Certifiée} : $0{,}4954739718\pm2\cdot10^{-13}$ ; noyau extrapolé $0{,}49547385$\\
Perte de $C^2$ de $P_0(t;\cdot)$, premier épisode (thm.~\ref{prop:premier}) & protocole $\Pi$ & Démontrée (assistée par ordinateur) ; coefficient $c_\varepsilon$ à $\varepsilon$ fixé\\
$T_1\notin C^2$ (prop.~\ref{cor:T1}) & protocole $\Pi$ & Démontrée (inversion écrite) ; pente $-2{,}3867$ vs $-2{,}3872\cdot10^{-2}$ par décade\\
Contact $a_\varepsilon\ge0$ & --- & Pas de mécanisme structurel trouvé\\
Persistance du mot près de $\mathrm{Fix}(R)$, réf. non axiales, $d\ge3$, profondeur & --- & Ouvert\\\bottomrule
\end{tabular}
\end{center}
\noindent En résumé, la dichotomie se lit ainsi : \emph{dans la classe isotrope, la régularité $C^2$ et la borne $O(\varepsilon)$ survivent aux zéros de $V_0$ par construction ; lorsque l'anisotropie historique $\lambda>0$ intervient et que la réponse transverse satisfait (N), l'imparité impose un zéro du déplacement sur le sous-espace axial, et (N) le rend régulier, d'où une perte de $C^2$ au voisinage de ce sous-espace.} Cette séparation est celle qu'on soupçonnait ; elle n'est pas encore un théorème général (conditions (N), persistance du mot).

\end{document}
